\section{探索性分析与 Worktree 隔离}

探索性数据分析（EDA）是 Codex 最受益于干净隔离的阶段。一个 worktree 可以测试地址清洗或特征工程，另一个专注于图表或替代建模方向。

\subsection{使用 Git Worktree 管理并行探索}

Codex App 内置了 worktree 支持。在终端中也可以直接使用 Git worktree：

\begin{lstlisting}[caption={创建 worktree 隔离不同探索方向},language=bash]
git worktree add ../analysis-highway-eda \
    -b analysis/highway-eda
git worktree add ../analysis-model-comparison \
    -b analysis/highway-modeling
\end{lstlisting}

\begin{knowledgebox}{为什么要用 Worktree？}
Worktree 的价值在于：
\begin{itemize}[nosep]
    \item 每个 diff 可独立 review，不会混入无关改动
    \item 防止一个长线程中混入不兼容的想法
    \item 保持每条探索路径的可回溯性
\end{itemize}
\end{knowledgebox}

在贯穿示例中，这一步要做的是：比较靠近高速公路的房屋与远离高速公路的房屋，检查异常值，审视缺失值模式，判断观察到的效应是真实的还是反映了社区构成、房屋面积或其他混杂因素。

\subsection{本章小结}
EDA 阶段利用 Worktree 将不同假设和探索方向物理隔离，保持每个分支的 diff 清晰可审查。

